\documentclass[11pt]{article}
\usepackage[margin=1in]{geometry}
\usepackage{amsmath,amssymb}
\usepackage{booktabs}
\usepackage{graphicx}
\usepackage{hyperref}
\usepackage{natbib}
\usepackage{caption}
\usepackage{subcaption}

\title{Survival Analysis of ChatGPT Conversation Duration:\\
Do Follow-Up Suggestions Extend Engagement?}
\author{AI Disclosures Project}
\date{January 2026}

\begin{document}

\maketitle

\begin{abstract}
We analyze conversation duration in ChatGPT using a discrete-time survival framework, examining whether follow-up suggestions extend conversations or whether observed increases in conversation length are driven by model capability improvements. Using 91,343 turn-level observations from 42,692 conversations (2023--2025), we find that (1) the tail of conversation length is getting heavier over time (Pareto $\alpha$ decreasing from 3.51 to 2.82), (2) suggestions are associated with \textit{higher} hazard of conversation ending (OR = 1.14, $p < 0.001$), suggesting a selection effect rather than a causal engagement mechanism, and (3) model capability and year effects are highly collinear, making it difficult to disentangle their contributions.
\end{abstract}

\section{Introduction}

Large language model interfaces increasingly incorporate ``engagement features'' such as suggested follow-up prompts, designed to encourage users to continue conversations. A key question for AI governance and product design is whether these features genuinely extend engagement or merely appear at natural conversation endpoints.

This analysis applies survival analysis methods to model the probability of conversation continuation conditional on turn number, model version, year, and whether suggestions were offered.

\section{Data}

\subsection{Source}

We use data from WildChat, a dataset of ChatGPT conversations collected via browser extensions. The data includes turn-level observations with:

\begin{itemize}
    \item \textbf{Conversation identifiers}: Unique hash per conversation
    \item \textbf{Turn structure}: Turn number within conversation, total turns
    \item \textbf{Continuation outcome}: Whether user sent another message
    \item \textbf{Suggestion indicators}: Count of suggestions shown, whether any were used
    \item \textbf{Temporal variables}: Timestamp, year (2023--2025)
    \item \textbf{Model information}: GPT version (e.g., gpt-4o-mini, gpt-3.5-turbo)
    \item \textbf{Content features}: User and assistant word counts
\end{itemize}

\subsection{Sample Characteristics}

\begin{table}[h]
\centering
\caption{Conversation Distribution by Year}
\label{tab:desc_stats}
\begin{tabular}{lrrrrrrr}
\toprule
Year & $n$ & Mean & SD & Median & P90 & P95 & P99 \\
\midrule
2023 & 14,915 & 2.32 & 2.73 & 1 & 5 & 7 & 14 \\
2024 & 13,040 & 1.93 & 2.81 & 1 & 4 & 6 & 15 \\
2025 & 14,737 & 2.14 & 3.42 & 1 & 4 & 7 & 18 \\
\bottomrule
\end{tabular}
\end{table}

A key feature of the data is the \textbf{heavy right tail}: while the median conversation length is 1 turn across all years, the 99th percentile ranges from 14 to 18 turns, with maximum lengths of 42--50 turns.

\section{Methods}

\subsection{Discrete-Time Survival Framework}

We model conversation duration using a discrete-time hazard framework. Let $T$ denote the number of turns in a conversation. At each turn $t$, we observe whether the conversation continues ($Y_t = 0$) or ends ($Y_t = 1$), conditional on survival to turn $t$.

The discrete-time hazard is:
\begin{equation}
h(t) = P(T = t \mid T \geq t) = P(Y_t = 1 \mid Y_1 = 0, \ldots, Y_{t-1} = 0)
\end{equation}

We model this using logistic regression:
\begin{equation}
\text{logit}[h(t \mid \mathbf{X})] = \beta_0 + \beta_1 \log(t) + \beta_2 \cdot \text{suggestion} + \beta_3 \cdot \text{year}_{2024} + \beta_4 \cdot \text{year}_{2025}
\end{equation}

The log-transform of turn number allows a flexible baseline hazard that captures the declining hazard observed empirically.

\subsection{Power-Law Tail Analysis}

To characterize the heavy tail, we fit a Pareto distribution to the survival function for $t \geq 5$:
\begin{equation}
P(T \geq t) \propto t^{-\alpha}
\end{equation}

where $\alpha$ is the tail exponent. Lower $\alpha$ indicates a heavier tail (more mass in extreme values).

\subsection{Split-Population (Cure) Model}

Given that 68.7\% of conversations end at turn 1, we consider a split-population model that separates:
\begin{enumerate}
    \item $\pi$: Probability of being a ``multi-turn'' user
    \item $S(t \mid \text{uncured})$: Survival function for those who might continue
\end{enumerate}

This is implemented by fitting separate models for first-turn and later-turn hazards.

\subsection{Identification Challenge}

A key methodological challenge is that GPT model versions are highly collinear with year:
\begin{itemize}
    \item 2023: Dominated by gpt-3.5-turbo-0613, gpt-3.5-turbo-0301, gpt-4-0314
    \item 2024: Dominated by gpt-4o-2024-08-06, gpt-4o-mini-2024-07-18
    \item 2025: Dominated by gpt-4o-mini-2024-07-18, gpt-4.1-mini-2025-04-14
\end{itemize}

Only gpt-4o-mini-2024-07-18 spans two years (2024--2025), providing limited within-model variation for separating year effects from model capability effects.

\section{Results}

\subsection{Tail Is Getting Heavier Over Time}

\begin{table}[h]
\centering
\caption{Pareto Tail Exponent by Year}
\label{tab:pareto}
\begin{tabular}{lrrr}
\toprule
Year & Pareto $\alpha$ & $R^2$ & $P(T > 10)$ \\
\midrule
2023 & 3.51 & 0.97 & 2.3\% \\
2024 & 3.18 & 0.89 & 2.1\% \\
2025 & 2.82 & 0.90 & 2.7\% \\
\bottomrule
\end{tabular}
\end{table}

The Pareto exponent $\alpha$ is \textbf{decreasing} over time, indicating that the tail of conversation length is getting heavier. More conversations are reaching extreme lengths in 2025 compared to 2023.

\subsection{Hazard Declines with Turn Number}

\begin{table}[h]
\centering
\caption{Empirical Hazard by Turn Number}
\label{tab:hazard}
\begin{tabular}{lrrr}
\toprule
Turn & $n$ at risk & Hazard & Continuation Rate \\
\midrule
1 & 42,692 & 68.7\% & 31.3\% \\
2 & 13,350 & 38.4\% & 61.6\% \\
3 & 8,220 & 30.9\% & 69.1\% \\
5 & 4,108 & 23.5\% & 76.5\% \\
10 & 1,284 & 20.3\% & 79.7\% \\
15 & 487 & 14.2\% & 85.8\% \\
\bottomrule
\end{tabular}
\end{table}

The hazard of ending declines sharply with turn number, consistent with positive duration dependence: users who survive early turns are increasingly likely to continue.

\subsection{Discrete-Time Hazard Model}

\begin{table}[h]
\centering
\caption{Baseline Discrete-Time Hazard Model (Logistic Regression)}
\label{tab:baseline}
\begin{tabular}{lrrrr}
\toprule
Variable & $\beta$ & SE & OR & 95\% CI \\
\midrule
Intercept & 0.385 & 0.017 & 1.469 & [1.42, 1.52] \\
$\log(\text{turn})$ & $-$1.078 & 0.010 & 0.340 & [0.33, 0.35] \\
Suggestion given & 0.135 & 0.016 & 1.145 & [1.11, 1.18] \\
Year 2024 & 0.343 & 0.018 & 1.410 & [1.36, 1.46] \\
Year 2025 & 0.238 & 0.017 & 1.269 & [1.23, 1.31] \\
\bottomrule
\end{tabular}
\begin{flushleft}
\small{Note: $N = 91,343$ turn-level observations. All coefficients significant at $p < 0.001$. Pseudo $R^2 = 0.136$.}
\end{flushleft}
\end{table}

\textbf{Key findings:}
\begin{enumerate}
    \item \textbf{Suggestion effect is \textit{positive} (counterintuitive)}: OR = 1.145 means suggestions are associated with 14.5\% \textit{higher} odds of ending. This suggests a \textbf{selection effect}: suggestions appear at natural stopping points, not that they cause earlier termination.

    \item \textbf{Year effects}: Both 2024 (OR = 1.41) and 2025 (OR = 1.27) show higher hazard than 2023, after controlling for turn number and suggestions.

    \item \textbf{Turn effect}: OR = 0.34 for $\log(\text{turn})$ confirms strong negative duration dependence---the longer a conversation has lasted, the lower the hazard of ending.
\end{enumerate}

\subsection{Split-Population Analysis}

Separating first-turn from later-turn behavior reveals different dynamics:

\begin{table}[h]
\centering
\caption{First-Turn vs. Later-Turn Models}
\label{tab:split}
\begin{tabular}{lrrrr}
\toprule
& \multicolumn{2}{c}{First Turn} & \multicolumn{2}{c}{Later Turns} \\
\cmidrule(lr){2-3} \cmidrule(lr){4-5}
Variable & OR & 95\% CI & OR & 95\% CI \\
\midrule
Suggestion given & 1.098 & [1.05, 1.15] & 1.198 & [1.15, 1.25] \\
Year 2024 & 1.381 & [1.31, 1.46] & 1.449 & [1.37, 1.54] \\
Year 2025 & 1.152 & [1.10, 1.21] & 1.410 & [1.34, 1.49] \\
\bottomrule
\end{tabular}
\end{table}

The suggestion effect is stronger in later turns (OR = 1.20) than at first turn (OR = 1.10), consistent with suggestions appearing at more critical decision points as conversations progress.

\section{Discussion}

\subsection{Interpreting the Suggestion Effect}

The finding that suggestions are associated with \textit{higher} hazard of ending requires careful interpretation. Three explanations are possible:

\begin{enumerate}
    \item \textbf{Selection effect}: The algorithm shows suggestions precisely when it predicts the user is about to leave. The association is not causal but reflects the algorithm's targeting.

    \item \textbf{Suggestions as ``closing'' signals}: Suggestions may inadvertently signal ``here are some things you could ask, but we're done with the current topic,'' prompting users to end.

    \item \textbf{Suggestions improve quality at cost of length}: If suggestions help users get what they need faster, this could increase satisfaction while reducing conversation length.
\end{enumerate}

Without randomization of suggestion display, we cannot distinguish these mechanisms.

\subsection{Model Capability vs. Engagement Design}

The heavy collinearity between model version and year prevents clean identification of whether:
\begin{itemize}
    \item Newer models are more capable, resolving queries faster (shorter conversations)
    \item Newer models are more engaging, leading to longer conversations
    \item Platform design changes (suggestions, UI) drive the changes
\end{itemize}

The fact that 2024 shows the \textit{highest} hazard (shortest conversations) despite more advanced models suggests capability effects may dominate: better models resolve queries faster.

\subsection{The Growing Heavy Tail}

The decreasing Pareto $\alpha$ (3.51 → 2.82) indicates more extreme-length conversations over time. Possible explanations:
\begin{enumerate}
    \item \textbf{Power users}: A growing segment uses ChatGPT for extended tasks (coding, writing, analysis)
    \item \textbf{Model capability}: Better models support more complex, multi-turn interactions
    \item \textbf{Use case expansion}: New applications (agentic workflows, research assistance) inherently require more turns
\end{enumerate}

\section{Conclusion}

This analysis provides three key findings:

\begin{enumerate}
    \item The tail of conversation duration is getting \textbf{heavier over time}, with more extreme-length conversations in 2025.

    \item Follow-up suggestions are associated with \textbf{higher hazard of ending}, suggesting they appear at natural stopping points rather than extending engagement.

    \item \textbf{Model capability and year effects cannot be cleanly separated}, requiring either instrumental variables or randomized experiments for causal identification.
\end{enumerate}

These findings have implications for AI platform design and transparency: reported engagement metrics should distinguish between ``more conversations'' and ``longer conversations,'' and the role of suggestion algorithms in shaping user behavior warrants further investigation.

\section*{Data and Code Availability}

Data and analysis code available at: \url{https://github.com/SrulyRosenblat/follow_up_research}

\end{document}
